% ═══════════════════════════════════════════════════════════════
% ML-03: Elektrische Leistung & Energie - MUSTERLÖSUNG
% Physik Klasse 8 | Greenhouse Schools Rostock
% ═══════════════════════════════════════════════════════════════

\documentclass[11pt, a4paper]{article}

\usepackage[utf8]{inputenc}
\usepackage[T1]{fontenc}
\usepackage[ngerman]{babel}

\usepackage{helvet}
\renewcommand{\familydefault}{\sfdefault}

\usepackage{microtype}
\usepackage[margin=1.5cm, top=2cm, bottom=1.5cm]{geometry}
\usepackage{enumitem}
\usepackage{tabularx}
\usepackage{booktabs}
\usepackage{array}
\usepackage{multicol}
\usepackage{tikz}
\usetikzlibrary{calc}

\usepackage{amsmath, amssymb}
\usepackage{siunitx}
\sisetup{locale = DE, per-mode = fraction, fraction-function = \dfrac}

\usepackage{fancyhdr}
\usepackage{lastpage}
\usepackage{needspace}

\usepackage{xcolor}
\usepackage{tcolorbox}
\tcbuselibrary{skins}

% ═══════════════════════════════════════════════════════════════
% FARBEN
% ═══════════════════════════════════════════════════════════════

\definecolor{physik}{HTML}{2c5aa0}
\definecolor{hellgrau}{HTML}{F5F5F5}
\definecolor{mittelgrau}{HTML}{888888}
\definecolor{rahmen}{HTML}{CCCCCC}
\definecolor{loesung}{HTML}{C62828}

% ═══════════════════════════════════════════════════════════════
% VARIABLEN
% ═══════════════════════════════════════════════════════════════

\newcommand{\fach}{Physik}
\newcommand{\thema}{Elektrische Leistung \& Energie}
\newcommand{\klasse}{8}
\newcommand{\stunde}{03}
\newcommand{\maxpunkte}{25}

\newif\ifzeigesternchen
\zeigesternchentrue

\colorlet{fachfarbe}{physik}

% ═══════════════════════════════════════════════════════════════
% KOPF- UND FUSSZEILE
% ═══════════════════════════════════════════════════════════════

\pagestyle{fancy}
\fancyhf{}
\renewcommand{\headrulewidth}{0.4pt}
\renewcommand{\footrulewidth}{0pt}
\lhead{\small \fach{} Klasse \klasse{} | \thema{} | \textcolor{loesung}{\textbf{MUSTERLÖSUNG}}}
\rhead{}
\cfoot{\small\textcolor{mittelgrau}{Seite \thepage{} von \pageref{LastPage}}}

% ═══════════════════════════════════════════════════════════════
% BOXEN
% ═══════════════════════════════════════════════════════════════

\newtcolorbox{merkkasten}[1][]{
    colback=hellgrau,
    colframe=hellgrau,
    boxrule=0pt,
    arc=0pt,
    left=8pt, right=8pt, top=8pt, bottom=6pt,
    borderline north={2pt}{0pt}{fachfarbe},
    before skip=6pt, after skip=6pt,
    #1
}

% ═══════════════════════════════════════════════════════════════
% BEFEHLE
% ═══════════════════════════════════════════════════════════════

\newcommand{\lsg}[1]{\textcolor{loesung}{\textbf{#1}}}
\newcommand{\punktebox}[1]{\hfill\small\textcolor{mittelgrau}{[#1]}}
\newcommand{\stern}[1]{\ifzeigesternchen\textsuperscript{\textcolor{fachfarbe}{(#1)}}\fi}

\newcommand{\aufgabe}[2]{%
    \needspace{4\baselineskip}%
    \vspace{10pt}%
    \noindent\textbf{\textcolor{fachfarbe}{Aufgabe #1}} #2%
    \vspace{4pt}%
    \nopagebreak[4]%
}

\newcommand{\operator}[1]{\textbf{\textcolor{fachfarbe}{#1}}}

\newlist{teil}{enumerate}{2}
\setlist[teil,1]{label=\alph*), leftmargin=1.8em, itemsep=3pt, parsep=0pt, topsep=3pt}

\newcommand{\ein}[1]{\,\mathrm{#1}}

\setlength{\parindent}{0pt}
\setlength{\parskip}{4pt}

% ═══════════════════════════════════════════════════════════════
% DOKUMENT
% ═══════════════════════════════════════════════════════════════

\begin{document}

% HEADER
\begin{center}
    {\Large\bfseries\textcolor{fachfarbe}{\thema}}\\[0.2em]
    {\small Arbeitsblatt | \fach{} Klasse \klasse{} | Stunde \stunde{} | \textcolor{loesung}{\textbf{MUSTERLÖSUNG}}}
\end{center}
\vspace{-0.3em}
\hrule
\vspace{0.6em}

% ═══════════════════════════════════════════════════════════════
% MERKKASTEN 1: Leistung
% ═══════════════════════════════════════════════════════════════

\begin{merkkasten}
\textbf{Elektrische Leistung}

Die elektrische \lsg{Leistung} gibt an, wie viel \lsg{Energie} pro \lsg{Sekunde} umgewandelt wird.

\vspace{0.3em}
\textbf{Definition:} \quad $P = \dfrac{E}{t}$ \quad (Leistung = Energie geteilt durch Zeit)

\vspace{0.3em}
\textbf{Berechnung:} \quad $P = $ \lsg{U $\cdot$ I} \qquad
\textbf{Einheit:} \quad $[P] = 1$ \lsg{Watt} $= 1\ein{W}$

\vspace{0.3em}
\textbf{Umrechnung:} \quad $1\ein{kW} = $ \lsg{1000} $\ein{W}$
\end{merkkasten}

% ═══════════════════════════════════════════════════════════════
% MERKKASTEN 2: Energie
% ═══════════════════════════════════════════════════════════════

\begin{merkkasten}
\textbf{Elektrische Energie}

Elektrische Energie ist \lsg{Leistung} mal \lsg{Zeit}. \quad (Umstellung: $E = P \cdot t$)

\vspace{0.3em}
\textbf{Formel:} \quad $E = $ \lsg{P $\cdot$ t}

\vspace{0.3em}
\textbf{SI-Einheit:} \quad $1\ein{Ws} = 1\ein{J}$ \quad (Joule ist die SI-Einheit der Energie)

\vspace{0.3em}
\textbf{Haushaltsüblich:} \quad $1\ein{kWh} = $ \lsg{3.600.000} $\ein{J}$ \quad (für Stromabrechnung)

\vspace{0.3em}
\textbf{Umrechnung:} \quad $1\ein{kWh} = $ \lsg{1000} $\ein{Wh}$
\end{merkkasten}

% ═══════════════════════════════════════════════════════════════
% AUFGABE 1: Geräteetiketten
% ═══════════════════════════════════════════════════════════════

\aufgabe{1}{\stern{*}} \punktebox{6}

\operator{Suche} im Physikraum drei elektrische Geräte mit Typenschild. \operator{Notiere} die Angaben.

\vspace{0.3em}
\renewcommand{\arraystretch}{1.6}
\begin{tabularx}{\linewidth}{|p{3.5cm}|X|X|X|X|}
\hline
\textbf{Gerät} & \textbf{Spannung (V)} & \textbf{Stromstärke (A)} & \textbf{Leistung (W)} & \textbf{Sonstiges} \\
\hline
1. \lsg{z.B. Netzgerät} & \lsg{230} & \lsg{0,5} & \lsg{–} & \lsg{50 Hz} \\
\hline
2. \lsg{z.B. iPad-Ladegerät} & \lsg{5} & \lsg{2,4} & \lsg{12} & \lsg{USB} \\
\hline
3. \lsg{z.B. Neonröhre} & \lsg{230} & \lsg{–} & \lsg{36} & \lsg{T8} \\
\hline
\end{tabularx}
\renewcommand{\arraystretch}{1.0}

\textcolor{loesung}{\small Hinweis: Individuelle Lösungen möglich. Wichtig: Korrekte Einheiten erkannt.}

% ═══════════════════════════════════════════════════════════════
% AUFGABE 2: Berechnen
% ═══════════════════════════════════════════════════════════════

\aufgabe{2}{\stern{*}} \punktebox{4}

\operator{Berechne} die elektrische Leistung.

\begin{teil}
    \item Ein Gerät wird mit $U = 12\ein{V}$ und $I = 3\ein{A}$ betrieben.

    \textcolor{loesung}{\small $P = U \cdot I = 12\ein{V} \cdot 3\ein{A} = 36\ein{W}$}
    \hfill $P = $ \lsg{36 W}

    \item Ein Föhn zieht bei $U = 230\ein{V}$ einen Strom von $I = 10\ein{A}$.

    \textcolor{loesung}{\small $P = U \cdot I = 230\ein{V} \cdot 10\ein{A} = 2300\ein{W}$}
    \hfill $P = $ \lsg{2300 W}
\end{teil}

% ═══════════════════════════════════════════════════════════════
% AUFGABE 3: Energie und kWh
% ═══════════════════════════════════════════════════════════════

\aufgabe{3}{\stern{**}} \punktebox{6}

\operator{Berechne} die elektrische Energie in Wh und kWh.

\begin{teil}
    \item Eine Lampe mit $P = 60\ein{W}$ leuchtet $t = 5\ein{h}$.

    \vspace{2em}
    $E = $ \lsg{P $\cdot$ t = 60 W $\cdot$ 5 h} $= $ \lsg{300} Wh $= $ \lsg{0,3} kWh

    \item Ein Computer mit $P = 200\ein{W}$ läuft $t = 4\ein{h}$ pro Tag, 30 Tage im Monat.

    \vspace{2.5em}
    $E_{\text{Monat}} = $ \lsg{200 W $\cdot$ 4 h $\cdot$ 30 = 24.000 Wh} $= $ \lsg{24} kWh
\end{teil}

% ═══════════════════════════════════════════════════════════════
% AUFGABE 4: Stromkosten
% ═══════════════════════════════════════════════════════════════

\aufgabe{4}{\stern{**}} \punktebox{5}

\operator{Berechne} die Stromkosten. Der Strompreis beträgt $\mathbf{30\,\dfrac{ct}{kWh}}$.

\vspace{0.3em}
Ein Durchlauferhitzer hat $P = 20\ein{kW}$. Familie Müller duscht täglich insgesamt 15 Minuten.

\begin{teil}
    \item Wie viel Energie wird pro Tag verbraucht?

    \textcolor{loesung}{\small (15 Min = 0,25 h)}
    \vspace{1em}
    $E = $ \lsg{20 kW $\cdot$ 0,25 h} $= $ \lsg{5} kWh

    \item Wie viel kostet das Duschen im Monat (30 Tage)?

    \vspace{2em}
    Kosten $= $ \lsg{5 kWh $\cdot$ 30 Tage $\cdot$ 30 $\dfrac{\text{ct}}{\text{kWh}}$ = 4500 ct} $= $ \lsg{45} Euro
\end{teil}

% ═══════════════════════════════════════════════════════════════
% AUFGABE 5: Transfer
% ═══════════════════════════════════════════════════════════════

\aufgabe{5}{\stern{***}} \punktebox{4}

\operator{Erkläre}, warum Energieversorger die Einheit \textbf{kWh} statt \textbf{Joule} verwenden.

\vspace{0.5em}
\textcolor{loesung}{
\textbf{Musterantwort:}\\
1 kWh = 3.600.000 J. Die Zahl in Joule wäre viel zu groß und unhandlich für Abrechnungen. Mit kWh erhält man übersichtliche Zahlen (z.B. 3000 kWh/Jahr statt 10.800.000.000 J). Außerdem entspricht kWh der praktischen Nutzung: Geräte haben Leistungsangaben in Watt/kW und laufen Stunden, nicht Sekunden.
}

% ═══════════════════════════════════════════════════════════════
% ERWEITERUNG: Haushaltsanalyse
% ═══════════════════════════════════════════════════════════════

\vspace{1em}
\hrule
\vspace{0.8em}

\noindent\textbf{\textcolor{fachfarbe}{Erweiterung (Hausaufgabe):}} \hfill \textit{freiwillig}

\vspace{0.3em}
\noindent\operator{Analysiere} den Stromverbrauch in deinem Haushalt. \operator{Recherchiere} die Leistungsaufnahme typischer Geräte und \operator{berechne} den Jahresverbrauch.

\vspace{0.3em}
\small
\textit{\textbf{Hinweis:} Bei Geräten wie Kühlschrank oder Waschmaschine steht auf dem Typenschild die \textbf{maximale} Leistung. Der tatsächliche Verbrauch ist geringer, da diese Geräte nicht dauerhaft mit voller Leistung laufen. Recherchiere typische \textbf{Durchschnittswerte} (z.B. Kühlschrank $\approx$ 100--150 kWh/Jahr, Waschmaschine $\approx$ 0,5--1 kWh pro Waschgang).}
\normalsize

\vspace{0.5em}
\renewcommand{\arraystretch}{1.8}
\begin{tabularx}{\linewidth}{|p{3.2cm}|X|X|X|X|}
\hline
\textbf{Gerät} & \textbf{Leistung} & \textbf{Laufzeit/Tag} & \textbf{Energie/Jahr} & \textbf{Kosten/Jahr} \\
 & (W oder kW) & (h) & (kWh) & (€, bei 30 ct/kWh) \\
\hline
Kühlschrank & \lsg{--} & \lsg{durchgehend} & \lsg{120} & \lsg{36} \\
\hline
Fernseher & \lsg{100 W} & \lsg{4 h} & \lsg{146} & \lsg{44} \\
\hline
Computer/Laptop & \lsg{60 W} & \lsg{3 h} & \lsg{66} & \lsg{20} \\
\hline
Smartphone-Laden & \lsg{10 W} & \lsg{2 h} & \lsg{7} & \lsg{2} \\
\hline
Licht (gesamt) & \lsg{50 W} & \lsg{5 h} & \lsg{91} & \lsg{27} \\
\hline
Waschmaschine & \lsg{--} & \lsg{3×/Woche} & \lsg{78} & \lsg{23} \\
\hline
\lsg{Trockner} & \lsg{--} & \lsg{2×/Woche} & \lsg{200} & \lsg{60} \\
\hline
\lsg{Geschirrspüler} & \lsg{--} & \lsg{täglich} & \lsg{250} & \lsg{75} \\
\hline
\multicolumn{3}{|r|}{\textbf{Summe:}} & \lsg{$\approx$ 960} & \lsg{$\approx$ 290} \\
\hline
\end{tabularx}
\renewcommand{\arraystretch}{1.0}

\vspace{0.3em}
\textcolor{loesung}{\small\textit{Beispielwerte – individuelle Lösungen variieren je nach Haushalt!}}

\vspace{0.5em}
\textbf{a)} Welche drei Geräte verbrauchen in deinem Haushalt am meisten Strom?

\vspace{0.3em}
\textcolor{loesung}{Typische Großverbraucher: \textbf{1.} Geschirrspüler/Trockner, \textbf{2.} Kühlschrank (läuft 24/7), \textbf{3.} Fernseher oder Computer. Auch Warmwasserbereitung (Durchlauferhitzer, Boiler) sind große Verbraucher!}

\vspace{0.5em}
\textbf{b)} \operator{Schlage} zwei konkrete Maßnahmen vor, um Strom zu sparen.

\vspace{0.3em}
\textcolor{loesung}{\textbf{Mögliche Maßnahmen:}\\
• Standby vermeiden (Steckerleisten mit Schalter)\\
• Wäsche bei niedrigerer Temperatur waschen (30°C statt 60°C)\\
• LED-Lampen statt Glühbirnen verwenden\\
• Kühlschrank nicht neben Heizung/Herd stellen\\
• Geräte mit guter Energieeffizienzklasse (A+++) kaufen\\
• Nur volle Wasch-/Spülmaschinen laufen lassen}

% ═══════════════════════════════════════════════════════════════
% PUNKTE
% ═══════════════════════════════════════════════════════════════

\vfill
\hrule
\vspace{0.4em}

\begin{tabularx}{\linewidth}{X r}
\small\textcolor{mittelgrau}{$\star$ Grundlagen: 10 BE | $\star\star$ Anwendung: 11 BE | $\star\star\star$ Transfer: 4 BE} &
\textbf{Punkte:} \maxpunkte{} / \maxpunkte
\end{tabularx}

\end{document}
